\documentclass[11pt,a4paper]{article}
\usepackage[margin=2.5cm]{geometry}
\usepackage[T1]{fontenc}
\usepackage{lmodern}
\usepackage{xcolor}
\usepackage{hyperref}
\begin{document}

\begin{center}
{\Large\bfseries INTERNAL DRAFT — HOLD FOR COMPARATOR GATE}
\end{center}

\noindent\textbf{Date:} To be completed at submission

\vspace{1em}
\noindent\textbf{To:} The Editors \\
\textit{International Journal of Production Economics}

\vspace{1em}
\noindent\textbf{Subject:} Revised manuscript on perishable inventory-routing under demand uncertainty

\vspace{1em}
Dear Editors,

This manuscript examines replenishment, age-specific inventory carryover, expiry, customer service, and shared fleet capacity in a multi-period perishable inventory-routing setting. It treats decision-focused learning as one candidate approach and evaluates it alongside forecast-based and policy-based alternatives.

The revised study uses a hand-specified synthetic demand process with independently generated episodes. It does not use or calibrate to M5 observations and contains no operational stock, expiry, service, or routing records. All reported schedules are evaluated by a common FIFO simulator, and solver limits, feasible construction fallbacks, and non-optimal statuses are disclosed. We make no claim of field validation or state-of-the-art performance.

The related-work section acknowledges direct stochastic and robust perishable inventory-routing studies and the recent scenario predict-then-optimize work for online inventory routing. A post-hoc finite-scenario robust adaptation has been evaluated, but it is not a reproduction of the published robust methods and all its candidate MILP calls use fallback. The manuscript must not be submitted until a credible source-method comparator gate is closed and all claims and results are synchronized to that completed evidence.

\vspace{1em}
Sincerely,

\vspace{2em}
The authors

\end{document}

